\chapter{สถาปัตยกรรมการวิเคราะห์คุณภาพดินระดับดิจิทัลและตัวแปรทางชีวกายภาพ}
\label{chap:ch05}

\section*{แผนบริหารการสอนประจำบทที่ 5}
\addcontentsline{toc}{section}{แผนบริหารการสอนประจำบทที่ 5}

\noindent\textbf{หัวข้อเนื้อหาประจำบท}
\begin{enumerate}
    \item คุณสมบัติทางกายภาพและเนื้อดินในสวนทุเรียนภาคตะวันออก
    \item คุณสมบัติทางเคมี ดัชนีความอุดมสมบูรณ์ และค่า pH ดิน
    \item การเปรียบเทียบระหว่างวิธีวิเคราะห์ดินในห้องแล็บกับเซนเซอร์ดิจิทัล
\end{enumerate}

\noindent\textbf{วัตถุประสงค์เชิงพฤติกรรม}
\begin{enumerate}
    \item อธิบายหลักการ ทฤษฎี และสถาปัตยกรรมเทคโนโลยีประจำบทที่ 5 ได้อย่างถูกต้อง
    \item ประยุกต์ใช้เครื่องมือ อุปกรณ์เซนเซอร์ และระบบปัญญาประดิษฐ์เพื่อการจัดการสวนทุเรียนได้อย่างมีประสิทธิภาพ
    \item วิเคราะห์ สรุปผลข้อมูล และแก้ไขปัญหาเชิงฟิสิกส์เกษตรและคุณภาพดินได้อย่างเป็นระบบ
\end{enumerate}

\noindent\textbf{การวัดและประเมินผล}
\begin{table}[H]
\centering\small
\begin{tabularx}{\textwidth}{p{0.55\textwidth} p{0.25\textwidth} c}
\toprule
\textbf{เกณฑ์การประเมินผล} & \textbf{เครื่องมือวัดผล} & \textbf{สัดส่วนคะแนน} \\
\midrule
1. ทักษะการปฏิบัติงานและการใช้อุปกรณ์ & แบบประเมินทักษะภาคสนาม & 40\% \\
2. รายงานข้อมูลและการวิเคราะห์ผล & การตรวจสมุดบันทึกและดาต้าชีท & 30\% \\
3. แบบฝึกหัดและคำถามท้ายบท & แบบทดสอบท้ายบทเรียน & 20\% \\
4. การมีส่วนร่วมและความรับผิดชอบ & การเข้าเรียนและการตรงต่อเวลา & 10\% \\
\midrule
\textbf{รวมทั้งสิ้น} & & \textbf{100\%} \\
\bottomrule
\end{tabularx}
\end{table}
\newpage

% ==========================================================
% เนื้อหาบทเรียนประจำบทที่ 5
% ==========================================================

\section{คุณสมบัติทางกายภาพและเนื้อดินในสวนทุเรียนภาคตะวันออก}
\index{คุณสมบัติทางกายภาพและเนื้อดินในสวนทุเรียนภาคตะวันออก}

การศึกษาและทำความเข้าใจเกี่ยวกับคุณสมบัติทางกายภาพและเนื้อดินในสวนทุเรียนภาคตะวันออก (Soil Physical Properties and Texture) มีความสำคัญอย่างยิ่งในระบบเกษตรอัจฉริยะ โดยมีรายละเอียดสาระสำคัญดังนี้

ดินในพื้นที่ปลูกทุเรียนแถบจังหวัดจันทบุรีและระยองส่วนใหญ่จัดอยู่ในกลุ่มดินชุดท่าใหม่ (Tha Mai series) และชุดสระแก้ว ซึ่งพัฒนามาจากหินบะซอลต์และหินตะกอน มีลักษณะเป็นดินร่วนเหนียวปนทรายหรือดินร่วนปนดินเหนียวสีแดง

คุณสมบัติทางกายภาพวิกฤตของดินทุเรียน:
- ความหนาแน่นรวมของดิน (Bulk density, $\rho_b$): ดินที่ดีควรมีค่า $\rho_b$ อยู่ในช่วง 1.1 ถึง 1.3 g/cm$^3$ หากดินแน่นทึบเกิน 1.5 g/cm$^3$ จะขัดขวางการแทงทะลุของรากแขนงและการถ่ายเทอากาศ
- ความพรุนรวมของดิน (Total Soil Porosity): ควรมีสัดส่วนประมาณ 50\% แบ่งเป็นช่องพรุนขนาดใหญ่ (Macropores สำหรับระบายน้ำและถ่ายเทอากาศ) 25\% และช่องพรุนขนาดเล็ก (Micropores สำหรับกักเก็บน้ำ) 25\%
- อัตราการแทรกซึมของน้ำ (Infiltration Rate): ทุเรียนต้องการดินที่มีการระบายน้ำดีเยี่ยม น้ำต้องไม่ขังแฉะเกิน 24 ชั่วโมง

\section{คุณสมบัติทางเคมี ดัชนีความอุดมสมบูรณ์ และค่า pH ดิน}
\index{คุณสมบัติทางเคมี ดัชนีความอุดมสมบูรณ์ และค่า pH ดิน}

การศึกษาและทำความเข้าใจเกี่ยวกับคุณสมบัติทางเคมี ดัชนีความอุดมสมบูรณ์ และค่า pH ดิน (Soil Chemical Properties, pH and Fertility) มีความสำคัญอย่างยิ่งในระบบเกษตรอัจฉริยะ โดยมีรายละเอียดสาระสำคัญดังนี้

ทุเรียนเจริญเติบโตได้ดีที่สุดในดินที่มีค่าความเป็นกรด-ด่าง (pH) อยู่ในช่วง 5.5 ถึง 6.5

ผลกระทบของค่า pH ดินต่อความเป็นประโยชน์ของธาตุอาหาร:
1. ดินกรดรุนแรง (pH $< 5.0$): ฟอสฟอรัสจะถูกตรึงด้วยเหล็กและอะลูมิเนียม เกิดพิษจากอะลูมิเนียมและแมงกานีส ส่งผลให้ปลายรากทุเรียนกุดและหยุดเจริญเติบโต
2. ดินด่าง (pH $> 7.5$): จุลธาตุอาหาร เช่น เหล็ก สังกะสี และทองแดง จะตกตะกอน ทำให้ทุเรียนแสดงอาการใบเหลืองซีด
3. ค่าการนำไฟฟ้าของดิน (Soil EC): ค่า EC ของสารละลายอิ่มตัวในดินไม่ควรเกิน 1.5 dS/m เพื่อป้องกันแรงดันออสโมซิสดึงน้ำออกจากรากพืช

\section{การเปรียบเทียบระหว่างวิธีวิเคราะห์ดินในห้องแล็บกับเซนเซอร์ดิจิทัล}
\index{การเปรียบเทียบระหว่างวิธีวิเคราะห์ดินในห้องแล็บกับเซนเซอร์ดิจิทัล}

การศึกษาและทำความเข้าใจเกี่ยวกับการเปรียบเทียบระหว่างวิธีวิเคราะห์ดินในห้องแล็บกับเซนเซอร์ดิจิทัล (Laboratory Soil Standards vs Digital Probes) มีความสำคัญอย่างยิ่งในระบบเกษตรอัจฉริยะ โดยมีรายละเอียดสาระสำคัญดังนี้

การวิเคราะห์คุณภาพดินตามมาตรฐานสากลและห้องปฏิบัติการปฐพีวิทยา:
- ค่า pH ดิน: วัดด้วยวิธีสารแขวนลอยดินต่อน้ำกลั่น 1:1 หรือ 1:2.5 ด้วย pH Meter หัววัดกระจก
- ปริมาณอินทรียวัตถุ (Organic Matter; OM): วิเคราะห์ด้วยวิธี Walkley-Black Wet Oxidation
- ไนโตรเจนทั้งหมด (Total N): วิธี Semi-micro Kjeldahl
- ฟอสฟอรัสที่เป็นประโยชน์ (Available P): สกัดด้วยวิธี Bray II
- โพแทสเซียมที่แลกเปลี่ยนได้ (Exchangeable K): สกัดด้วย 1N Ammonium acetate และวัดด้วย Flame Photometer

เซนเซอร์ดิจิทัลและโพรบภาคสนามช่วยให้เก็บข้อมูลได้อย่างต่อเนื่องและรวดเร็ว แต่ต้องผ่านกระบวนการสอบเทียบความถูกต้อง (Calibration) เทียบกับผลแล็บมาตรฐานเสมอเพื่อลดค่าความคลาดเคลื่อนเชิงระบบ (Systematic error)

\section{ขั้นตอนและวิธีปฏิบัติการทดลอง}
\index{ขั้นตอนการทดลองประจำบทที่ 5}

\subsection{การเก็บตัวอย่างดินสวนทุเรียนตามหลักปฐพีวิทยาเพื่อส่งวิเคราะห์แล็บ}
\begin{sensorbox}[การเก็บตัวอย่างดินสวนทุเรียนตามหลักปฐพีวิทยาเพื่อส่งวิเคราะห์แล็บ]
1. กำหนดจุดเก็บตัวอย่างแบบซิกแซก (Zig-zag) หรือตารางกริด (Grid sampling) 5-10 จุดต่อแปลง
2. กวาดเศษซากใบไม้แห้งและหญ้าบนผิวดินออก ใช้พลั่วขุดดินเป็นรูปตัว V ลึก 0-15 ซม. (ดินบน) และ 15-30 ซม. (ดินล่าง)
3. ปาดดินด้านข้างออก เหลือแถบดินตรงกลางกว้าง 1 นิ้ว ใส่ในถังพลาสติกสะอาด
4. คลุกเคล้าดินย่อยทั้งหมดให้เข้ากันเป็นเนื้อเดียว แบ่งดินตัวอย่างรวมหนัก 1 กิโลกรัม ใส่ถุงพลาสติก
5. ผึ่งลมในที่ร่มให้แห้ง บดผ่านตะแกรงขนาด 2 มิลลิเมตร พร้อมส่งวิเคราะห์
\end{sensorbox}

\subsection{การตรวจวัดค่า pH และ EC ดินภาคสนามด้วยเครื่องวัดดิจิทัลแบบพกพา}
\begin{sensorbox}[การตรวจวัดค่า pH และ EC ดินภาคสนามด้วยเครื่องวัดดิจิทัลแบบพกพา]
1. ชั่งดินแห้งที่ร่อนผ่านตะแกรง 20 กรัม ใส่ลงในบีกเกอร์
2. เติมน้ำกลั่นบริสุทธิ์ 20 มิลลิลิตร (อัตราส่วนดินต่อน้ำ 1:1)
3. คนให้เข้ากันอย่างต่อเนื่องเป็นเวลา 30 นาที เพื่อให้เกิดสมดุลเคมี
4. สอบเทียบหัววัด pH ด้วยสารละลายบัฟเฟอร์มาตรฐาน pH 4.01 และ 7.00
5. จุ่มหัววัดลงในส่วนของสารแขวนลอย อ่านค่า pH และบันทึกผลเมื่อตัวเลขนิ่ง
\end{sensorbox}

\subsection{การประเมินเนื้อสัมผัสของดินภาคสนามด้วยวิธีสัมผัส}
\begin{sensorbox}[การประเมินเนื้อสัมผัสของดินภาคสนามด้วยวิธีสัมผัส]
1. หยิบดินตัวอย่างประมาณ 1 ช้อนโต๊ะ พรมน้ำให้อยู่ในสภาพดินชื้นพอปั้นได้
2. คลึงดินบนฝ่ามือเป็นแท่งกลมเส้นผ่านศูนย์กลางประมาณ 3-5 มิลลิเมตร
3. ดันดินเป็นเส้นยาว (Ribbon test): หากทำเส้นยาวได้เกิน 5 ซม. แสดงว่ามีดินเหนียวสูง
4. ขยี้ดินเปียกบนฝ่ามือเพื่อสัมผัสความสาก (ทราย), ความนุ่มลื่น (ทรายละเอียด/ทรายแป้ง), หรือความเหนียวเหนอะหนะ (ดินเหนียว)
\end{sensorbox}

\section{ตารางบันทึกผลการทดลองและการอภิปรายผล}
\begin{table}[H]
\centering\small
\caption{ตารางบันทึกผลการตรวจวัดและข้อมูลพารามิเตอร์ประจำบทที่ 5}
\label{tab:ch05_datasheet}
\begin{tabularx}{\textwidth}{p{0.3\textwidth} p{0.3\textwidth} X}
\toprule
\textbf{พารามิเตอร์ / การทดสอบ} & \textbf{ค่าที่ตรวจวัดได้} & \textbf{การแปลผลและการวิเคราะห์} \\
\midrule
จุดตรวจวัดที่ 1 (บริเวณดินดอน) & ค่าความชื้นอยู่ในเกณฑ์ปกติ & ปริมาณน้ำเพียงพอต่อการเจริญ \\
จุดตรวจวัดที่ 2 (บริเวณที่ลุ่ม) & มีการสะสมความชื้นสูง & ควรชะลอการให้น้ำเพื่อป้องกันรากเน่า \\
สถานะสัญญาณโครงข่าย & RSSI: -85 dBm, SNR: +8 dB & คุณภาพสัญญาณ LoRaWAN ยอดเยี่ยม \\
\bottomrule
\end{tabularx}
\end{table}

\section{คำถามทบทวนและแบบฝึกหัดท้ายบท}
\begin{enumerate}
    \item จงอธิบายความสัมพันธ์ระหว่างปริมาณอินทรียวัตถุในดิน (OM) กับความจุในการกักเก็บน้ำของดินชุดท่าใหม่
    \item เหตุใดการใส่ปุ๋ยเคมีปริมาณมากเกินไปจึงส่งผลให้ค่า EC ของดินพุ่งสูง และส่งผลเสียต่อรากทุเรียนอย่างไร
    \item จงเปรียบเทียบข้อดีและข้อจำกัดของการวิเคราะห์ดินด้วยเครื่องแล็บมาตรฐานกับการใช้เซนเซอร์ดิจิทัลวัดในแปลง
\end{enumerate}
